%% GENERATED by python/paper/build.py from results/ -- do not edit by hand.
%% Rebuild:  .venv\Scripts\python.exe python\paper\build.py --only ArgentinaUniversal
%% Source:   results/shocks/{eeOnly,universal}_match_rho1.0000_commonX.csv
\begin{table}[!htb]
\centering
\begin{threeparttable}
\caption{Pension system reform, year 2010.}
\label{table:Argentina:Universal}
\begin{tabular}{lccc}
\toprule
\textbf{Scenario} & \textbf{Tax rate} & \textbf{Savings rate} & \textbf{Avg. workweek (hours)} \\
\midrule 
Baseline & 10.92\% & 14.40\% & 42.54 \\[1ex]
Economic Equilibrium & 10.92\% & $+0.18$ p.p. & 42.50 \\[1ex]
Full effect & 12.17\% & $-0.28$ p.p. & 42.39 \\[1ex]
\bottomrule
\end{tabular}
\begin{tablenotes}
\footnotesize
\item \textit{Note:} The reform permanently shifts $\epsilon = 1-\theta + \theta \eta_1 h_1/h$ at the calibration year, unanticipated. The ``Economic Equilibrium'' scenario solves for the new economic equilibrium path with taxes as in the Baseline scenario, but with the new permanent $\epsilon$; the ``Full effect'' scenario lets taxes be determined by the politico-economic equilibrium. $\rho=1$. The savings rate is savings relative to GDP; the baseline row reports its level and every other row the change against that baseline in percentage points. Aggregate hours have no scale in the model, so the workweek is normalised to the observed average of 42.54 hours in the baseline. Common-$X$ calibration: see the note to Table~\ref{table:Arg:Calib}.
\end{tablenotes}
\end{threeparttable}
\end{table}
